\section{3/9}
\subsection{What is a program?}
\begin{definition}
  A \textbf{Program} is a description of something computable.
\end{definition}
We have the question of \textunderscore{Syntax} vs. \textunderscore{Semantics}.
\\
Often times we talk about programming with turing-complete systems but this is not always true. Take English for example.
\\
English is comprised of sentences. 
\begin{lstlisting}
Sentence := (NounPhrase, VerbPhrase)
  NounPhrase := (determiner option, adjective option, noun)
  VerbPhrase := Simple of (verb, np option, adverb)
    | Rec (verb, VerbPhrase)
\end{lstlisting}
We have the sentence "Alice speaks" where Alice : noun and speaks : verb.
\\
Another: "(a black cat){NP} (can catch mice){VP}"
\\\\
This can be represented as a syntax tree \url{https://www.google.com/url?sa=t&source=web&rct=j&url=https%3A%2F%2Fwww.nltk.org%2Fbook%2Fch08.html&ved=0CBYQjRxqFwoTCIiemKmyk5MDFQAAAAAdAAAAABAI&opi=89978449}
\\
OCaml is also a language and also can be evaluated using a syntax tree.

\subsection{Creating a Calculator Language}
To create a language, you should start with examples of elements in the language.
\\
We have Lang Calc. with "*", "+", and $n \in \mathbb{N}$. 
\\
This can be specified in an OCaml type.
\begin{lstlisting}
type expr =
  | Nat of int
  | Add of expr * expr
  | Mul of expr * expr
\end{lstlisting}
We might want to make a serializer, a way to take this AST and turn it into a string.
\begin{lstlisting}
let rec string_of_expr e =
  match e with
  | Nat i -> string_of_int i
  | Add (l, r) -> string_of_expr l ^ " + " ^ string_of_expr r
  | Mul (l, r) -> string_of_expr l ^ " * " ^ string_of_expr r
\end{lstlisting}
If we have the statement string_of_expr (Mul(Nat 3, Add(Nat 4, Nat 2))) we get "3 * 4 + 2", but this ignores the (4 + 2). We need a round-trip property: 
\begin{lstlisting}
let rec string_of_expr e =
  match e with
  | Nat i -> string_of_int i
  | Add (l,r) -> "(" ^ string_of_expr l ^ " + " ^ string_of_expr r ^ ")"
  | Mul (l,r) -> "(" ^ string_of_expr l ^ " * " ^ string_of_expr r ^ ")"
\end{lstlisting}
\newpage
To parse a string in this language we create the token type
\begin{lstlisting}
type token =
  | NAT of int
  | PLUS
  | TIMES
  | LPAREN
  | RPAREN
  | EOF
\end{lstlisting}
To specify our roundtrip property, $\forall e \in$ Calc, parse (string_of_expr e)